% !TEX spellcheck = en_US



% ~~~~~~~~~~~~~~~~~~~~~~~~~~~~~~~~~~~~~~~ V
% (NC) 2026 Haluk Bingol
% github.com/halukbingol/LaTeX-Templates
%
% Licensees may copy, distribute, display, and perform the work and make derivative works 
% and remixes based on it only for non-commercial purposes. 
% ~~~~~~~~~~~~~~~~~~~~~~~~~~~~~~~~~~~~~~~ A




% =====================================================================
%  Generic Stack and Queue -- The Concepts
%  ---------------------------------------------------------------------
%  Built on hblib/Template/hblibTemplate.tex with the libraries in ../../hblib/.
%  Programs and their real outputs are in codeoutput/:
%     codeoutput/<Name>.java   a program (classes without main() are only compiled)
%     codeoutput/<Name>.txt    what it printed
%     codeoutput/JShell/       a jshell session
% =====================================================================

\documentclass[aspectratio=169,10pt, compress]{beamer}

% ~~~~~~~~~~~~~~~~~~~~~~~~~~~~~~~~~~~~~~~~~~~~~~~~~~~~~~~~~~~~~~~~~ V hblib
	\usepackage{../../hblib/hblibPresentation} % lib presentation: theme, i/N, \hSection, algorithm2e, ...
	\usepackage{../../hblib/hblibCodeoutput} % lib code + output: \lstcode, \lstterminal, ...
	\usepackage{../../hblib/hblibDatastructures} % lib data structures: \hString, \hStack, \hQueue
% ~~~~~~~~~~~~~~~~~~~~~~~~~~~~~~~~~~~~~~~~~~~~~~~~~~~~~~~~~~~~~~~~~ A hblib


% xcolor, array (and the L{width} column), listings, tikz and algorithm2e come from the hblib libraries
\usepackage[T1]{fontenc}
\usepackage{lmodern}
\usepackage{booktabs}

\newcommand{\win}{\textbf{Windows}}
\newcommand{\mac}{\textbf{macOS}}

% listings: fixed character cells, so the columns of a program's output stay aligned
\lstset{columns=fixed}

% algorithm2e set-up (the package is loaded by hblibPresentation)
\SetAlCapFnt{\small}
\SetAlCapNameFnt{\small}
\SetKwProg{Fn}{function}{}{end}
\SetKw{KwError}{error}

% ---------------------------------------------------------------------
\title{Generic Stack and Queue}
\subtitle{Set, Stack and Queue in the Java API}
\author[Alpaydin]{Ruhan Alpaydin}

\institute[FCIS]{
    Faculty of Computer and Information Sciences\\
    Yeditepe University
}
\date{
	{\tiny \hbTimeStamp}
}

\begin{document}




% ~~~~~~~~~~~~~~~~~~~~~~~~~~~~~~~~~~~~~~ V frm
\begin{frame}[t,fragile]
  \titlepage
\end{frame}
% ~~~~~~~~~~~~~~~~~~~~~~~~~~~~~~~~~~~~~~ A frm




% ~~~~~~~~~~~~~~~~~~~~~~~~~~~~~~~~~~~~~~ V frm
\begin{frame}[t,fragile] \frametitle{Outline}
  \begin{columns}[T]
    \begin{column}{0.48\textwidth}
      \tableofcontents[sections={1-3}, hideallsubsections]
    \end{column}
    \begin{column}{0.48\textwidth}
      \tableofcontents[sections={4-6}, hideallsubsections]
    \end{column}
  \end{columns}
\end{frame}
% ~~~~~~~~~~~~~~~~~~~~~~~~~~~~~~~~~~~~~~ A frm

% ~~~~~~~~~~~~~~~~~~~~~~~~~~~~~~~~~~~~~~ V
\hSection[Idea]{Three ADTs in One Picture}
% ~~~~~~~~~~~~~~~~~~~~~~~~~~~~~~~~~~~~~~ A

% ~~~~~~~~~~~~~~~~~~~~~~~~~~~~~~~~~~~~~~ V
\subsection[Picture]{The picture}
% ~~~~~~~~~~~~~~~~~~~~~~~~~~~~~~~~~~~~~~ A




% ~~~~~~~~~~~~~~~~~~~~~~~~~~~~~~~~~~~~~~ V frm
\begin{frame}[t,fragile] \frametitle{Set, Stack, Queue}
  \begin{columns}[T]
    \begin{column}{0.32\textwidth}
      \begin{block}{Set}
        \centering\small
        \hCode{\{apple, banana, cherry\}}\\[0.3em]
        \emph{no duplicates}
      \end{block}
      \small \hRemark{Is this element in the set?}
    \end{column}
    \begin{column}{0.32\textwidth}
      \begin{block}{Stack (LIFO)}
        \begin{center}
        \begin{tikzpicture}\hStack{0}{5}{3}{A,B,C,D}[0.5][0.4]\end{tikzpicture}
        \end{center}
        \small \hCode{pop} $\to$ \textbf{D}
      \end{block}
      \small \hRemark{Undo, brackets, backtracking.}
    \end{column}
    \begin{column}{0.32\textwidth}
      \begin{block}{Queue (FIFO)}
        \begin{center}
        \begin{tikzpicture}\hQueueArray{6}{0}{4}{4}{A,B,C,D}[0.5][0.4]\end{tikzpicture}
        \end{center}
        \small \hCode{dequeue} $\to$ \textbf{A}
      \end{block}
      \small \hRemark{Printers, buffers, BFS.}
    \end{column}
  \end{columns}
  \vspace{0.5em}
  \small All three are \hRemark{generic types} in the Java API: pick the element type in angle brackets.
\end{frame}
% ~~~~~~~~~~~~~~~~~~~~~~~~~~~~~~~~~~~~~~ A frm




% ~~~~~~~~~~~~~~~~~~~~~~~~~~~~~~~~~~~~~~ V frm
\begin{frame}[t,fragile] \frametitle{What ``generic'' means}
  \begin{columns}[T]
    \begin{column}{0.5\textwidth}
      \begin{itemize}
        \item A \hRemark{generic type} is parameterized by the type of its elements.
        \item \hCode{Set<String>}, \hCode{Stack<Integer>}, \hCode{Queue<Person>}.
        \item The \emph{same class} works for any reference type.
      \end{itemize}
    \end{column}
    \begin{column}{0.46\textwidth}
      \begin{block}{Why generics?}
        \small
        \begin{itemize}
          \item Type-safe: a wrong type is caught at \textbf{compile time}.
          \item No casting: \hCode{names.get(0)} is already a \hCode{String}.
        \end{itemize}
      \end{block}
    \end{column}
  \end{columns}
  \vspace{0.3em}
  \begin{center}\footnotesize
  \begin{tabular}{@{}lll@{}}
    \toprule
    \textbf{ADT}   & \textbf{Rule}                 & \textbf{Java type} \\
    \midrule
    \textbf{Set}   & no duplicates, no fixed order & \hCode{java.util.Set} \\
    \textbf{Stack} & LIFO                          & \hCode{Stack}, \hCode{Deque} \\
    \textbf{Queue} & FIFO                          & \hCode{java.util.Queue} \\
    \bottomrule
  \end{tabular}
  \end{center}
\end{frame}
% ~~~~~~~~~~~~~~~~~~~~~~~~~~~~~~~~~~~~~~ A frm

% ~~~~~~~~~~~~~~~~~~~~~~~~~~~~~~~~~~~~~~ V
\hSection[Set]{Set in the Java API}
% ~~~~~~~~~~~~~~~~~~~~~~~~~~~~~~~~~~~~~~ A

% ~~~~~~~~~~~~~~~~~~~~~~~~~~~~~~~~~~~~~~ V
\subsection[Concept]{What is a Set}
% ~~~~~~~~~~~~~~~~~~~~~~~~~~~~~~~~~~~~~~ A




% ~~~~~~~~~~~~~~~~~~~~~~~~~~~~~~~~~~~~~~ V frm
\begin{frame}[t,fragile] \frametitle{What is a Set?}
  \begin{itemize}
    \item A \hRemark{Set} is a collection of \textbf{unique} elements:
          no duplicates, ever.
    \item Order is \emph{not} part of the ADT unless the implementation
          says otherwise.
    \item Uses: keywords, visited cells, tags, membership tests,
          deduplication.
  \end{itemize}
  \begin{block}{Properties}
    \begin{itemize}
      \item Does not allow duplicate elements.
      \item Does not maintain insertion order (for \hCode{HashSet}).
    \end{itemize}
  \end{block}
\end{frame}
% ~~~~~~~~~~~~~~~~~~~~~~~~~~~~~~~~~~~~~~ A frm

% ~~~~~~~~~~~~~~~~~~~~~~~~~~~~~~~~~~~~~~ V
\subsection[Classes]{Three implementations}
% ~~~~~~~~~~~~~~~~~~~~~~~~~~~~~~~~~~~~~~ A




% ~~~~~~~~~~~~~~~~~~~~~~~~~~~~~~~~~~~~~~ V frm
\begin{frame}[t,fragile] \frametitle{Common Set implementations}
  \begin{center}\small
  \begin{tabular}{L{0.20\textwidth}L{0.26\textwidth}L{0.16\textwidth}L{0.26\textwidth}}
    \toprule
    \textbf{Class} & \textbf{Ordering}             & \textbf{Time}   & \textbf{When to use} \\
    \midrule
    \hCode{HashSet}       & none                          & $O(1)$ avg  & default; fast lookup \\
    \hCode{TreeSet}       & sorted (natural/comparator)   & $O(\log n)$ & sorted iteration \\
    \hCode{LinkedHashSet} & insertion order               & $O(1)$ avg  & predictable order \\
    \bottomrule
  \end{tabular}
  \end{center}
  \vspace{0.3em}
  \begin{itemize}\small
    \item All three implement \hCode{java.util.Set}.
    \item \hCode{HashSet} uses \hCode{hashCode()} and \hCode{equals()};
          override both in custom classes.
    \item \hCode{TreeSet} uses \hCode{compareTo()} (or a \hCode{Comparator}).
  \end{itemize}
\end{frame}
% ~~~~~~~~~~~~~~~~~~~~~~~~~~~~~~~~~~~~~~ A frm

% ~~~~~~~~~~~~~~~~~~~~~~~~~~~~~~~~~~~~~~ V
\subsection[Methods]{Key methods}
% ~~~~~~~~~~~~~~~~~~~~~~~~~~~~~~~~~~~~~~ A




% ~~~~~~~~~~~~~~~~~~~~~~~~~~~~~~~~~~~~~~ V frm
\begin{frame}[t,fragile] \frametitle{Key methods of \texttt{Set}}
  \begin{center}\small
  \begin{tabular}{lL{0.6\textwidth}}
    \toprule
    \textbf{Method}              & \textbf{Meaning} \\
    \midrule
    \hCode{add(E item)}         & add; \hCode{false} if it was already present \\
    \hCode{remove(Object o)}    & remove; \hCode{true} if found \\
    \hCode{contains(Object o)}  & membership test \\
    \hCode{isEmpty()}           & is there no element? \\
    \hCode{size()}              & number of elements \\
    \bottomrule
  \end{tabular}
  \end{center}
  \vspace{0.3em}
  \begin{itemize}\small
    \item \hCode{add} returns \hCode{boolean}: a free ``was it new?'' test.
    \item \hCode{contains} is the workhorse: sets exist for this call.
  \end{itemize}
\end{frame}
% ~~~~~~~~~~~~~~~~~~~~~~~~~~~~~~~~~~~~~~ A frm

% ~~~~~~~~~~~~~~~~~~~~~~~~~~~~~~~~~~~~~~ V
\subsection[Example]{In Java}
% ~~~~~~~~~~~~~~~~~~~~~~~~~~~~~~~~~~~~~~ A




% ~~~~~~~~~~~~~~~~~~~~~~~~~~~~~~~~~~~~~~ V frm
\begin{frame}[t,fragile] \frametitle{Set: example}
  \begin{columns}[T]
    \begin{column}{0.5\textwidth}
      \lstcode[basicstyle=\ttfamily\tiny]{codeoutput/SetDemo.java}{Java}
    \end{column}
    \begin{column}{0.46\textwidth}
      \lstterminal[basicstyle=\ttfamily\scriptsize]{codeoutput/SetDemo.txt}
      \small
      \begin{itemize}
        \item The second \hCode{add("Apple")} returns \hCode{false}.
        \item \hCode{HashSet} does \emph{not} promise any order.
        \item \hCode{LinkedHashSet} gives a predictable one.
      \end{itemize}
    \end{column}
  \end{columns}
\end{frame}
% ~~~~~~~~~~~~~~~~~~~~~~~~~~~~~~~~~~~~~~ A frm

% ~~~~~~~~~~~~~~~~~~~~~~~~~~~~~~~~~~~~~~ V
\hSection[Stack]{Stack in the Java API}
% ~~~~~~~~~~~~~~~~~~~~~~~~~~~~~~~~~~~~~~ A

% ~~~~~~~~~~~~~~~~~~~~~~~~~~~~~~~~~~~~~~ V
\subsection[Concept]{LIFO}
% ~~~~~~~~~~~~~~~~~~~~~~~~~~~~~~~~~~~~~~ A




% ~~~~~~~~~~~~~~~~~~~~~~~~~~~~~~~~~~~~~~ V frm
\begin{frame}[t,fragile] \frametitle{Stack: last in, first out}
  \begin{columns}[T]
    \begin{column}{0.48\textwidth}
      \begin{itemize}
        \item Insert and remove at the \textbf{top}.
        \item The most recent element leaves first: \hRemark{LIFO}.
        \item Uses: undo, brackets, expression evaluation,
              backtracking, DFS.
      \end{itemize}
    \end{column}
    \begin{column}{0.48\textwidth}
      \begin{center}
      \begin{tikzpicture}\hStack{0}{6}{3}{A,B,C,D}[0.55][0.45]\end{tikzpicture}
      \end{center}
      \small \hCode{pop} returns \textbf{D}.
    \end{column}
  \end{columns}
\end{frame}
% ~~~~~~~~~~~~~~~~~~~~~~~~~~~~~~~~~~~~~~ A frm

% ~~~~~~~~~~~~~~~~~~~~~~~~~~~~~~~~~~~~~~ V
\subsection[Methods]{Class java.util.Stack}
% ~~~~~~~~~~~~~~~~~~~~~~~~~~~~~~~~~~~~~~ A




% ~~~~~~~~~~~~~~~~~~~~~~~~~~~~~~~~~~~~~~ V frm
\begin{frame}[t,fragile] \frametitle{Class: \texttt{java.util.Stack}}
  \begin{center}\small
  \begin{tabular}{lL{0.55\textwidth}}
    \toprule
    \textbf{Method} & \textbf{Meaning} \\
    \midrule
    \hCode{push(E item)}   & add on top \\
    \hCode{pop()}          & remove and return the top; error if empty \\
    \hCode{peek()}         & return the top without removing it \\
    \hCode{isEmpty()}      & is there no element? \\
    \bottomrule
  \end{tabular}
  \end{center}
  \vspace{0.3em}
  \begin{itemize}\small
    \item \hCode{Stack<E>} is a legacy class (Java 1.0) that extends
          \hCode{Vector}: every call is \hCode{synchronized}.
    \item Modern Java prefers \hCode{Deque<E>} (\hCode{ArrayDeque})
          with the same \hCode{push}/\hCode{pop}/\hCode{peek} names.
  \end{itemize}
\end{frame}
% ~~~~~~~~~~~~~~~~~~~~~~~~~~~~~~~~~~~~~~ A frm

% ~~~~~~~~~~~~~~~~~~~~~~~~~~~~~~~~~~~~~~ V
\subsection[Modern]{Modern alternative}
% ~~~~~~~~~~~~~~~~~~~~~~~~~~~~~~~~~~~~~~ A




% ~~~~~~~~~~~~~~~~~~~~~~~~~~~~~~~~~~~~~~ V frm
\begin{frame}[t,fragile] \frametitle{Modern alternative: \texttt{Deque}}
  \begin{columns}[T]
    \begin{column}{0.52\textwidth}
      \begin{itemize}
        \item Same three operations: \hCode{push}, \hCode{pop}, \hCode{peek}.
        \item \hCode{ArrayDeque} is faster than \hCode{Stack}, no legacy baggage.
        \item \hRemark{Rule of thumb}: \hCode{ArrayDeque} for both
              stack and queue in new code.
      \end{itemize}
    \end{column}
    \begin{column}{0.44\textwidth}
      \begin{block}{Two refusals}
        \small
        \begin{itemize}
          \item \hCode{ArrayDeque} forbids \hCode{null}.
          \item \hCode{pop} on empty throws
                \hCode{NoSuchElementException}.
        \end{itemize}
      \end{block}
    \end{column}
  \end{columns}
\end{frame}
% ~~~~~~~~~~~~~~~~~~~~~~~~~~~~~~~~~~~~~~ A frm

% ~~~~~~~~~~~~~~~~~~~~~~~~~~~~~~~~~~~~~~ V
\subsection[Example]{In Java}
% ~~~~~~~~~~~~~~~~~~~~~~~~~~~~~~~~~~~~~~ A




% ~~~~~~~~~~~~~~~~~~~~~~~~~~~~~~~~~~~~~~ V frm
\begin{frame}[t,fragile] \frametitle{Stack: example}
  \begin{columns}[T]
    \begin{column}{0.5\textwidth}
      \lstcode[basicstyle=\ttfamily\tiny]{codeoutput/StackDemo.java}{Java}
    \end{column}
    \begin{column}{0.46\textwidth}
      \lstterminal[basicstyle=\ttfamily\scriptsize]{codeoutput/StackDemo.txt}
      \small After the two \hCode{pop}s the stack is empty: a
      further \hCode{pop} would throw
      \hCode{EmptyStackException}.
    \end{column}
  \end{columns}
\end{frame}
% ~~~~~~~~~~~~~~~~~~~~~~~~~~~~~~~~~~~~~~ A frm

% ~~~~~~~~~~~~~~~~~~~~~~~~~~~~~~~~~~~~~~ V
\hSection[Queue]{Queue in the Java API}
% ~~~~~~~~~~~~~~~~~~~~~~~~~~~~~~~~~~~~~~ A

% ~~~~~~~~~~~~~~~~~~~~~~~~~~~~~~~~~~~~~~ V
\subsection[Concept]{FIFO}
% ~~~~~~~~~~~~~~~~~~~~~~~~~~~~~~~~~~~~~~ A




% ~~~~~~~~~~~~~~~~~~~~~~~~~~~~~~~~~~~~~~ V frm
\begin{frame}[t,fragile] \frametitle{Queue: first in, first out}
  \begin{columns}[T]
    \begin{column}{0.48\textwidth}
      \begin{itemize}
        \item Insert at the \textbf{rear}, remove at the \textbf{front}.
        \item The oldest element leaves first: \hRemark{FIFO}.
        \item Uses: printers, keyboard buffers, scheduling, BFS,
              message queues.
      \end{itemize}
    \end{column}
    \begin{column}{0.48\textwidth}
      \begin{center}
      \begin{tikzpicture}\hQueueArray{6}{0}{4}{4}{A,B,C,D}[0.55][0.45]\end{tikzpicture}
      \end{center}
      \small \hCode{dequeue} returns \textbf{A}.
    \end{column}
  \end{columns}
\end{frame}
% ~~~~~~~~~~~~~~~~~~~~~~~~~~~~~~~~~~~~~~ A frm

% ~~~~~~~~~~~~~~~~~~~~~~~~~~~~~~~~~~~~~~ V
\subsection[Interface]{The Queue interface}
% ~~~~~~~~~~~~~~~~~~~~~~~~~~~~~~~~~~~~~~ A




% ~~~~~~~~~~~~~~~~~~~~~~~~~~~~~~~~~~~~~~ V frm
\begin{frame}[t,fragile] \frametitle{\texttt{java.util.Queue}: two forms}
  \begin{itemize}
    \item \hCode{Queue} is an \textbf{interface}, not a class.
    \item Main implementations: \hCode{ArrayDeque}, \hCode{LinkedList}.
    \item Every operation has two forms: one \textbf{throws},
          one returns a \textbf{special value}.
  \end{itemize}
  \begin{center}\small
  \begin{tabular}{lll}
    \toprule
    \textbf{Operation} & \textbf{throws exception} & \textbf{returns special value} \\
    \midrule
    insert    & \hCode{add(e)}     & \hCode{offer(e)} \\
    remove    & \hCode{remove()}   & \hCode{poll()} \\
    examine   & \hCode{element()}  & \hCode{peek()} \\
    \bottomrule
  \end{tabular}
  \end{center}
  \small \hCode{isEmpty()} and \hCode{size()} complete the interface.
\end{frame}
% ~~~~~~~~~~~~~~~~~~~~~~~~~~~~~~~~~~~~~~ A frm

% ~~~~~~~~~~~~~~~~~~~~~~~~~~~~~~~~~~~~~~ V
\subsection[Example]{In Java}
% ~~~~~~~~~~~~~~~~~~~~~~~~~~~~~~~~~~~~~~ A




% ~~~~~~~~~~~~~~~~~~~~~~~~~~~~~~~~~~~~~~ V frm
\begin{frame}[t,fragile] \frametitle{Queue: example}
  \begin{columns}[T]
    \begin{column}{0.52\textwidth}
      \lstcode[basicstyle=\ttfamily\tiny]{codeoutput/QueueDemo.java}{Java}
    \end{column}
    \begin{column}{0.44\textwidth}
      \lstterminal[basicstyle=\ttfamily\scriptsize]{codeoutput/QueueDemo.txt}
      \small \hCode{ArrayDeque} prints from the \textbf{front};
      \hCode{poll()} on empty returns \hCode{null}.
    \end{column}
  \end{columns}
\end{frame}
% ~~~~~~~~~~~~~~~~~~~~~~~~~~~~~~~~~~~~~~ A frm

% ~~~~~~~~~~~~~~~~~~~~~~~~~~~~~~~~~~~~~~ V
\subsection[Classes]{Which class}
% ~~~~~~~~~~~~~~~~~~~~~~~~~~~~~~~~~~~~~~ A




% ~~~~~~~~~~~~~~~~~~~~~~~~~~~~~~~~~~~~~~ V frm
\begin{frame}[t,fragile] \frametitle{Which \texttt{Queue} class?}
  \begin{center}\footnotesize
  \begin{tabular}{@{}lL{0.16\textwidth}L{0.42\textwidth}@{}}
    \toprule
    \textbf{Class}               & \textbf{Since} & \textbf{Note} \\
    \midrule
    \hCode{ArrayDeque}           & Java 6         & \hRemark{recommended}: a growing circular array \\
    \hCode{LinkedList}           & Java 1.2       & a doubly linked list \\
    \hCode{PriorityQueue}        & Java 5         & \hRemark{not FIFO}: the smallest element leaves first \\
    \hCode{ArrayBlockingQueue}   & Java 5         & bounded; waits when full or empty \\
    \hCode{ConcurrentLinkedQueue}& Java 5         & unbounded, lock-free, for many threads \\
    \bottomrule
  \end{tabular}
  \end{center}
  \vspace{0.3em}
  \begin{itemize}\small
    \item \hCode{ArrayDeque} is faster and more compact than \hCode{LinkedList};
          it forbids \hCode{null}.
    \item \hCode{PriorityQueue} is a \hCode{Queue} by interface only, \emph{not} by order.
  \end{itemize}
\end{frame}
% ~~~~~~~~~~~~~~~~~~~~~~~~~~~~~~~~~~~~~~ A frm

% ~~~~~~~~~~~~~~~~~~~~~~~~~~~~~~~~~~~~~~ V
\hSection[Exercises]{Exercises}
% ~~~~~~~~~~~~~~~~~~~~~~~~~~~~~~~~~~~~~~ A

% ~~~~~~~~~~~~~~~~~~~~~~~~~~~~~~~~~~~~~~ V
\subsection[Set]{Set exercises}
% ~~~~~~~~~~~~~~~~~~~~~~~~~~~~~~~~~~~~~~ A




% ~~~~~~~~~~~~~~~~~~~~~~~~~~~~~~~~~~~~~~ V frm
\begin{frame}[t,fragile] \frametitle{Exercise 1: remove duplicates using a Set}
  \begin{block}{Problem}
    Write a Java program to remove duplicates from a list using a
    \hCode{Set}.
  \end{block}
  \begin{columns}[T]
    \begin{column}{0.52\textwidth}
      \lstcode[basicstyle=\ttfamily\tiny]{codeoutput/RemoveDup.java}{Java}
    \end{column}
    \begin{column}{0.44\textwidth}
      \lstterminal[basicstyle=\ttfamily\scriptsize]{codeoutput/RemoveDup.txt}
      \small
      \hCode{LinkedHashSet} keeps the order of first occurrence;
      \hCode{HashSet} is faster but loses the order.
    \end{column}
  \end{columns}
\end{frame}
% ~~~~~~~~~~~~~~~~~~~~~~~~~~~~~~~~~~~~~~ A frm




% ~~~~~~~~~~~~~~~~~~~~~~~~~~~~~~~~~~~~~~ V frm
\begin{frame}[t,fragile] \frametitle{Exercise 2: common elements in two lists}
  \begin{block}{Problem}
    Write a Java program to find common elements between two lists
    using a \hCode{Set}.
  \end{block}
  \begin{columns}[T]
    \begin{column}{0.52\textwidth}
      \lstcode[basicstyle=\ttfamily\tiny]{codeoutput/CommonElements.java}{Java}
    \end{column}
    \begin{column}{0.44\textwidth}
      \lstterminal[basicstyle=\ttfamily\scriptsize]{codeoutput/CommonElements.txt}
      \small
      \hCode{retainAll} keeps only the elements also present in the
      argument: the intersection of the two sets.
    \end{column}
  \end{columns}
\end{frame}
% ~~~~~~~~~~~~~~~~~~~~~~~~~~~~~~~~~~~~~~ A frm

% ~~~~~~~~~~~~~~~~~~~~~~~~~~~~~~~~~~~~~~ V
\subsection[Stack]{Stack exercise}
% ~~~~~~~~~~~~~~~~~~~~~~~~~~~~~~~~~~~~~~ A




% ~~~~~~~~~~~~~~~~~~~~~~~~~~~~~~~~~~~~~~ V frm
\begin{frame}[t,fragile] \frametitle{Exercise 3: reverse a string using a Stack}
  \begin{block}{Problem}
    Write a Java program to reverse a string using a \hCode{Stack}.
  \end{block}
  \begin{columns}[T]
    \begin{column}{0.52\textwidth}
      \lstcode[basicstyle=\ttfamily\tiny]{codeoutput/ReverseString.java}{Java}
    \end{column}
    \begin{column}{0.44\textwidth}
      \lstterminal[basicstyle=\ttfamily\scriptsize]{codeoutput/ReverseString.txt}
      \small
      Push every character, then pop them: the last pushed is the
      first to come out, so the string is reversed.
    \end{column}
  \end{columns}
\end{frame}
% ~~~~~~~~~~~~~~~~~~~~~~~~~~~~~~~~~~~~~~ A frm

% ~~~~~~~~~~~~~~~~~~~~~~~~~~~~~~~~~~~~~~ V
\subsection[Queue]{Queue exercise}
% ~~~~~~~~~~~~~~~~~~~~~~~~~~~~~~~~~~~~~~ A




% ~~~~~~~~~~~~~~~~~~~~~~~~~~~~~~~~~~~~~~ V frm
\begin{frame}[t,fragile] \frametitle{Exercise 4: queue from two stacks}
  \begin{block}{Problem}
    Write a Java program to implement a queue using two stacks.
  \end{block}
  \begin{columns}[T]
    \begin{column}{0.52\textwidth}
      \lstcode[basicstyle=\ttfamily\tiny]{codeoutput/QueueTwoStacks.java}{Java}
    \end{column}
    \begin{column}{0.44\textwidth}
      \lstterminal[basicstyle=\ttfamily\scriptsize]{codeoutput/QueueTwoStacks.txt}
      \small
      Two reversals compose to the original order: \hCode{in}
      receives, \hCode{out} serves; when \hCode{out} is empty,
      transfer everything.
    \end{column}
  \end{columns}
\end{frame}
% ~~~~~~~~~~~~~~~~~~~~~~~~~~~~~~~~~~~~~~ A frm

% ~~~~~~~~~~~~~~~~~~~~~~~~~~~~~~~~~~~~~~ V
\hSection[Summary]{Summary}
% ~~~~~~~~~~~~~~~~~~~~~~~~~~~~~~~~~~~~~~ A




% ~~~~~~~~~~~~~~~~~~~~~~~~~~~~~~~~~~~~~~ V frm
\begin{frame}[t,fragile] \frametitle{Big ideas}
  \begin{itemize}
    \item \textbf{Set}: unique elements. \hCode{HashSet} (unordered),
          \hCode{TreeSet} (sorted), \hCode{LinkedHashSet} (insertion order).
    \item \textbf{Stack}: LIFO. \hCode{java.util.Stack} is legacy;
          prefer \hCode{Deque} (\hCode{ArrayDeque}) with
          \hCode{push}/\hCode{pop}/\hCode{peek}.
    \item \textbf{Queue}: FIFO. Interface \hCode{Queue};
          \hCode{ArrayDeque} or \hCode{LinkedList}.
          \hCode{add}/\hCode{remove}/\hCode{element} throw,
          \hCode{offer}/\hCode{poll}/\hCode{peek} return special values.
    \item \hCode{PriorityQueue} implements \hCode{Queue}, but is
          \emph{not} FIFO.
    \item The API classes are \hRemark{generic}: one class, any element type.
  \end{itemize}
\end{frame}
% ~~~~~~~~~~~~~~~~~~~~~~~~~~~~~~~~~~~~~~ A frm




% ~~~~~~~~~~~~~~~~~~~~~~~~~~~~~~~~~~~~~~ V frm
\begin{frame}[t,fragile] \frametitle{Check yourself}
  \begin{enumerate}\small
    \item What does \hCode{Set.add} return? Why is that useful?
    \item \hCode{HashSet} vs \hCode{TreeSet} vs \hCode{LinkedHashSet}:
          what is the iteration order in each?
    \item \hCode{Stack.pop} on an empty stack: what happens?
          What about \hCode{Queue.poll}?
    \item Why is \hCode{ArrayDeque} preferred over \hCode{Stack} in new code?
    \item Which \hCode{Queue} methods throw? Which return a special value?
    \item Why is \hCode{PriorityQueue} not FIFO?
    \item Why does a queue made of two stacks behave like a FIFO?
  \end{enumerate}
\end{frame}
% ~~~~~~~~~~~~~~~~~~~~~~~~~~~~~~~~~~~~~~ A frm




% ~~~~~~~~~~~~~~~~~~~~~~~~~~~~~~~~~~~~~~ V frm
\begin{frame}[t,fragile] \frametitle{References}
  \begin{enumerate}
    \item Oracle Java Documentation.
    \item GeeksforGeeks.
    \item Joshua Bloch, \emph{Effective Java}
          (best practices for the Java collections framework).
    \item Herbert Schildt, \emph{Java: The Complete Reference}.
    \item \url{https://docs.oracle.com/javase/8/docs/api/}
    \item \url{https://docs.oracle.com/javase/tutorial/collections/}
  \end{enumerate}
\end{frame}
% ~~~~~~~~~~~~~~~~~~~~~~~~~~~~~~~~~~~~~~ A frm




% ~~~~~~~~~~~~~~~~~~~~~~~~~~~~~~~~~~~~~~ V frm
\begin{frame}[t,fragile]
	\begin{center}
		\vspace{4cm}
		\hRemark{\Huge Questions?}
	\end{center}
\end{frame}
% ~~~~~~~~~~~~~~~~~~~~~~~~~~~~~~~~~~~~~~ A frm




\end{document}